\section{Introduction}
A radiograph usually reaches the archive with its anatomical region stored in a DICOM tag or a free-text study description, and either can be wrong or missing. An early audit of 562 images found false \emph{Body Part Examined} entries in 15.3 \% of them (Güld et al., SPIE 2002). Deep networks can read the region from the pixels instead (Pham et al. 2021; Langer et al. 2024; Mitsuyama et al. 2025). Most of that work uses private data, reports accuracy alone and says little about how the splits and thresholds were chosen.
On the public UNIFESP dataset (Farina, Abdala and Kitamura) we found one academic paper, which reports 86.94 \% accuracy for the best of five CNNs (Krishnan and Padmavathi 2025). That study treated the 1,738 training images as a single-label 22-class problem, added 1,088 GAN-generated images to the 1,686 originals, and reports accuracy only; we could not find a description of a patient-level split or of a held-out protocol, and a Kaggle competition whose final private leaderboard is led by a score of 0.93378. Two evaluation shortcuts are known to inflate results in other fields: leakage between images of the same patient, and thresholds or checkpoints chosen on data that is later scored (Kapoor and Narayanan 2023; Roberts et al. 2021). We did not find a measurement of how large these effects are for body-part classification A search of PubMed, arXiv and general web sources in November 2026 found studies of leakage in other imaging tasks but none that measured it for radiograph body-part classification..
So instead of chasing the leaderboard we asked two questions. How much of the reported performance survives a leakage-controlled protocol? And what would a clinical router need besides accuracy? The paper contributes:
\begin{itemize}\item a reproducible, patient-level benchmark on public data with confidence intervals (Tables 1 and 2);\item a measurement of how much common shortcuts inflate scores (Tables 4 and 5);\item negative results: no ablation, fusion variant or ensembling step was significantly better, and a gated CNN-Transformer fusion collapsed onto one branch (Tables 4 and 6);\item results on MURA and FracAtlas, including a body-part benchmark on the official MURA split (Tables 7 to 9);\item calibration, risk-coverage and a label audit based on out-of-fold predictions, aimed at routing use (Section 5.7);\item a frozen foundation-model baseline with a label-efficiency curve (Section 5.6) and a Grad-CAM faithfulness check (Table 10).\end{itemize}

